Formula (*) allows an easy determination of the natural filtration of $U_p(L)$ (cf. 1.3.6). Put $U=U_p(L)$, let $U^+$ be the two-sided ideal generated by $L$, and let $U_0=C\cdot1_U$. As in 1.3.6, one defines a filtration of $U$ (by sub-$C$-coalgebras) by putting, for $n\ge1$:
\[
U_n=\{u\in U\mid\Delta_U(u)-u\otimes1\in U_{n-1}\otimes U^+\}.
\]
Formula (*) then implies that $U_n$ is the \emph{free $C$-module with basis the monomials $x^m$ such that $|m|\le n$}.

\numpara{4.3.1}With the notation of 4.3, let us determine the elements $\xi$ of $U=U_p(L)$ such that $\varepsilon_U(\xi)=1$ and $\Delta_U(\xi)=\xi\otimes\xi$. Every element $\xi$ of $U$ is written $\xi=\sum_{n\in P}\xi_n(x^n/n!)$, with $\xi_n\in C$. The condition $\varepsilon_U(\xi)=1$ implies the equality $\xi_0=1$, where $0$ denotes the element of $P$ all of whose components are zero. The condition $\Delta_{U_p(L)}(\xi)=\xi\otimes\xi$ implies:
\[
\sum_{m\le n}\xi_n\frac{x^m}{m!}\otimes\frac{x^{n-m}}{(n-m)!}
=\sum_{q,r}\xi_q\xi_r\frac{x^q}{q!}\otimes\frac{x^r}{r!},
\]
that is,
\[
\xi_{q+r}=\xi_q\xi_r\quad\text{if }q_i+r_i<p,\qquad\xi_q\xi_r=0\quad\text{otherwise}.
\]
If one denotes by $\xi_i$ the value of $\xi_n$ when $n_i=1$ and $n_j=0$ for $j\ne i$, these conditions mean that one has
\[
\xi_n=\prod_i\xi_i^{n_i}\quad\text{if }n\in P,\qquad\xi_i^p=0\quad\forall i\in I.
\]
One derives from this:
\begin{proposition}{}
Let $k$ be a local pseudocompact ring\nde{We have added the hypothesis that $k$ be local, so that every topologically flat pseudocompact $k$-module is topologically free, cf. the proof.} of characteristic $p>0$, $\mathbf L$ a Lie $p$-algebra flat over $\mathbf O_k$, $C$ an object of $\mathbf{Alf}/k$ and $\sqrt[p]{0_C}$ the ideal of $C$ consisting of the elements $x$ such that $x^p=0$. There exists a bijection ``functorial in $C$'':
\[
\mathbf L(C)\otimes_C\sqrt[p]{0_C}\isomto\mathscr G_p(\mathbf L)(C).
\]
\end{proposition}
By Remark 1.2.3.F, one can indeed choose a basis $({}^Cx_i)_{i\in I}$ of $\mathbf L(C)$ over $C$ in such a way that, for every morphism $\varphi:C\to D$ of $\mathbf{Alf}/k$, $\mathbf L(\varphi)$ maps ${}^Cx_i$ to ${}^Dx_i$. By the preceding, the map
\[
\sum_i{}^Cx_i\otimes\xi_i\longmapsto\sum_{n\in P}\left(\prod_i\xi_i^{n_i}\right)\frac{{}^Cx^n}{n!}
\]
is a functorial bijection from $\mathbf L(C)\otimes_C\sqrt[p]{0_C}$ onto $\mathscr G_p(\mathbf L)(C)$.

\numpara{4.3.2}``\emph{There is no Campbell--Hausdorff formula in characteristic $p>0$}.'' Let us explain: the functorial isomorphism of 4.3.1 depends on the choice of the bases $({}^Cx_i)_{i\in I}$. Unlike what happens in 3.2 (the characteristic $0$ case), there is, in general, no bijection from $\mathbf L(C)\otimes_C\sqrt[p]{0_C}$ onto $\mathscr G_p(\mathbf L)(C)$ which is ``functorial in both $C$ and $\mathbf L$''. The argument which follows shows, for example, that such an isomorphism does not exist when $k$ is a field containing the $(p^2-1)$th roots of unity.

Indeed, choose $\mathbf L$ in such a way that, for every $C\in\mathbf{Alf}/k$, $\mathbf L(C)$ is the abelian Lie $p$-algebra over $C$ generated by an element $x$ subject to the relation $x^{(p^2)}=0$. One therefore has
\[
\mathbf L(C)=Cx\oplus Cx^{(p)},\qquad U_p(\mathbf L(C))\simeq k[x]/(x^{p^2}).
\]
% typo?: k[x], rather than C[x], is retained as printed.
We shall show that the only functorial morphism
\[
\chi_C:\mathbf L(C)\otimes_C\sqrt[p]{0_C}\longrightarrow U_p(\mathbf L(C))
\]
which is compatible with the endomorphisms of $\mathbf L$ and maps $\mathbf L(C)\otimes_C\sqrt[p]{0_C}$ into $\mathscr G_p(\mathbf L)(C)$ is the constant morphism with value $1$.

Indeed, let $\psi_C$ be the bijection from $\mathbf L(C)\otimes_C\sqrt[p]{0_C}$ onto $\mathscr G_p(\mathbf L)(C)=\operatorname{Prim}U_p(\mathbf L(C))$ given by 4.3.1, \ie:
\[
x\otimes a+x^{(p)}\otimes b\longmapsto\sum_{0\le i,j<p}a^ib^jx^{i+pj}.
\]
% typo?: Prim and the omitted factorial denominators are retained as printed.
Composing $\chi_C$ with $\psi_C^{-1}$, one obtains a morphism functorial (in $C$):
\[
\begin{aligned}
\theta_C:\quad\mathbf L(C)\otimes_C\sqrt[p]{0_C}&\longrightarrow\mathbf L(C)\otimes_C\sqrt[p]{0_C}\\
x\otimes a+x^{(p)}\otimes b&\longmapsto x\otimes P(a,b)+x^{(p)}\otimes Q(a,b).
\end{aligned}
\]
Functoriality in $C$ implies that $P(a,b)$ and $Q(a,b)$ are the values at $(a,b)$ of two polynomials $P,Q\in k[x,y]$ which one can suppose to have degree $<p$ in $X$ as well as in $Y$.\nde{We have detailed the original in what follows.} Let $\alpha$ be an element of $k$ and $\ell_\alpha$ the endomorphism of Lie $p$-algebra of $\mathbf L$ which maps $x$ to $\alpha x$ (and hence $x^{(p)}$ to $\alpha^px^{(p)}$). Then $U_p(\ell_\alpha)$ is the algebra endomorphism which sends $x$ to $\alpha x$ (and hence each $x^i$ to $\alpha^ix^i$, for $i<p^2$), and one sees easily that the square below
\[
\xymatrix@C=50pt@R=28pt{
\mathbf L(C)\otimes_C\sqrt[p]{0_C}\ar[r]^{\psi_C}\ar[d]_{\ell_\alpha(C)\otimes_C\id}&U_p(\mathbf L(C))\ar[d]^{U_p(\ell_\alpha)(C)}\\
\mathbf L(C)\otimes_C\sqrt[p]{0_C}\ar[r]_{\psi_C}&U_p(\mathbf L(C))}
\]
commutes. The hypothesis $\chi_C\circ(\ell_\alpha(C)\otimes\id)=U_p(\ell_\alpha)(C)\circ\chi_C$ then implies the equalities:
\[
P(\alpha a,\alpha^pb)=\alpha P(a,b)\quad\text{and}\quad Q(\alpha a,\alpha^pb)=\alpha^pQ(a,b).
\]
Writing $P=\sum_{i,j<p}\lambda_{ij}X^iY^j$ and $Q=\sum_{i,j<p}\mu_{ij}X^iY^j$, and taking for $C$ the algebra $k[X,Y]/(X^p,Y^p)$, one deduces that if $\lambda_{ij}\ne0$ (resp. $\mu_{ij}\ne0$) then $\alpha^{i-1+pj}=1$ (resp. $\alpha^{i+p(j-1)}=1$). Thus, taking for $\alpha$ a primitive root of unity of order $p^2-1$, one deduces that $P=\lambda X$ and $Q=\mu Y$, with $\lambda,\mu\in k$. Consequently, one has:
\[
\chi_C(x\otimes a+x^{(p)}\otimes b)=\sum_{0\le i,j<p}(\lambda a)^i(\mu b)^jx^{i+pj}.
\]
\footnotemark[138]Finally, consider the endomorphism $f$ of $\mathbf L$ which sends $x$ to $x+x^{(p)}$; taking $C=k[a]/(a^3)$ and comparing the coefficients of $x^p$ and $x^{p+1}$ in $(\chi_C\circ f(C))(x\otimes a)$ and in $(U_p(f)(C)\circ\chi_C)(x\otimes a)$, one obtains that $\lambda=\mu$ and $\lambda^2a^2=0$, whence $\lambda=\mu=0$.
% typo?: a^3 also for p=2 and the missing tensor identity on f(C) are retained.

\begin{theorem}{4.4}\label{VIIB.4.4}
Let $k$ be a pseudocompact ring of characteristic $p>0$ and $G$ a formal $k$-group. The following assertions are equivalent:
\begin{enumeratei}
\item There exists a Lie $p$-algebra $\mathbf L$ flat over $\mathbf O_k$ such that $G$ is isomorphic to $\mathscr G_p(\mathbf L)$ (and in this case $\mathbf L=\uLie(G)$ by 4.2.2).
\item There exists a projective pseudocompact $k$-module $\omega$ such that the affine algebra of $G$ is isomorphic to the quotient of $k[[\omega]]$ by the closed ideal generated by the $x^p$, for $x\in\omega$ (and in this case $\omega\simeq\omega_{G/k}$).
\item $G$ has height $\le1$ and $\omega_{G/k}$ is a projective pseudocompact $k$-module.
\item $G$ has height $\le1$ and is topologically flat over $k$.
\end{enumeratei}
\end{theorem}
\begin{proof}
(i) $\Rightarrow$ (ii): Let $\omega=\Gamma^*(\mathbf L)$ (cf. 1.2.3.D) and let $A$ be the quotient of $k[[\omega]]$ by the closed ideal generated by the $x^p$, for $x\in\omega$. Every continuous morphism $h:A\to C$, where $C$ is an object of $\mathbf{Alf}/k$, is determined by its restriction $h'$ to $\omega$; this restriction $h'$ sends $\omega$ into $\sqrt[p]{0_C}$. One thus obtains a canonical bijection from $\Hom_{\mathbf{Alp}/k}(A,C)$ onto the set $\Hom_c(\omega,\sqrt[p]{0_C})$ of continuous $k$-linear maps from $\omega$ to $\sqrt[p]{0_C}$. This last set is canonically isomorphic to $\mathbf L(C)\otimes_C\sqrt[p]{0_C}$ (see the proof of 3.3). The implication (i) $\Rightarrow$ (ii) therefore follows from the functorial bijection $\mathbf L(C)\otimes_C\sqrt[p]{0_C}\isomto\mathscr G_p(\mathbf L)(C)$ established in 4.3.1.

For the other implications, consult the proofs of Theorems 3.3 and VII A 7.4.1, which are analogous.
\end{proof}
\begin{remark}{4.4.A}\label{VIIB.4.4.A}\nde{We have added this remark, used in 4.4.2.}
Let $G$ be an infinitesimal formal $k$-group, with affine algebra $A$, such that $\omega_{G/k}=I_A/\overline{I_A^2}$ is a projective pseudocompact $k$-module. Then there exists a section $\tau:\omega_{G/k}\to I_A$ of the projection $I_A\to\omega_{G/k}$, and $\tau$ induces a continuous morphism of algebras $\psi:k[[\omega_{G/k}]]\to A$ which is \emph{surjective}, cf. the proof of the implication (iii) $\Rightarrow$ (i) in 3.3.
\end{remark}
\begin{corollary}{4.4.1}\label{VIIB.4.4.1}
If $k$ is a field of characteristic $p>0$, the functor $L\mapsto\mathscr G_p(L)$ is an equivalence from the category of Lie $p$-algebras over $k$ onto that of formal $k$-groups of height $\le1$.
\end{corollary}
Indeed, when $G$ ranges over the formal groups of height $\le1$, the functor $G\mapsto\mathscr G_p(\Lie G)$ is isomorphic to the identity functor by Theorem 4.4; similarly, we have seen in 4.2.2 that the functor $L\mapsto\Lie\mathscr G_p(L)$ is isomorphic to the identity functor (see also VII A, 5.5.3).\nde{If $k$ is an artinian ring of characteristic $p>0$, the same proof gives an equivalence between the category of flat Lie $p$-algebras over $k$ and that of topologically flat formal $k$-groups of height $\le1$.}

\numpara{4.4.2}Still take for $k$ a field of characteristic $p$. Let $G$ be an infinitesimal formal $k$-group, with affine algebra $A$.
% typo: un k-groupe formel infinitesimal G repeats G -> one occurrence.
Since $G$ is infinitesimal, every open ideal of $A$ contains $I_A^{\{p^n\}}$ for $n$ large enough, so $G$ is the projective limit of the affine algebras $A/I_A^{\{p^n\}}$ of the groups ${}_{\Fr^n}G$ (cf. 4.1.2). \emph{Every infinitesimal formal $k$-group is therefore an inductive limit of formal $k$-groups of finite height}.
% typo?: G, rather than A, is printed as the projective limit of affine algebras, retained.

Suppose that $G$ has height $\le n$ and put $H=G/{}_{\Fr}G$.\nde{We have detailed the original in what follows.} By 2.4 and 2.4.1, $\Fr:G\to G^{(p)}$ factors into an epimorphism $\pi:G\to H$ followed by a monomorphism $i:H\to G^{(p)}$. One then has the following commutative diagram:
\[
\xymatrix@C=24pt@R=28pt{
G\ar[r]^\pi\ar[dr]_{\Fr}&H\ar[r]^{\Fr^{n-1}}\ar[d]_i&H^{(p^{n-1})}\ar[d]^{i^{(p^{n-1})}}\\
&G^{(p)}\ar[r]_{\Fr^{n-1}}&G^{(p^n)}}
\]
and since the functor $X\mapsto X^{(p)}$ commutes with fibered products, $i^{(p^{n-1})}$ is again a monomorphism. Since $\Fr^n(G/k)$ is zero and $\pi$ is an epimorphism, $\Fr^{n-1}(G^{(p)}/k)\circ i=i^{(p^{n-1})}\circ\Fr^{n-1}(H/k)$ is zero and hence, since $i^{(p^{n-1})}$ is a monomorphism, $\Fr^{n-1}(H/k)$ is zero, so $H$ has height $\le n-1$. One therefore sees that: \emph{every formal $k$-group of finite height possesses a composition series whose quotients have height $\le1$}, hence can be described by Lie $p$-algebras over $k$.

Finally, the affine algebra $A$ of $G$ is a quotient of $k[[\omega_{G/k}]]$, cf. 4.4.A, so if $\omega_{G/k}$ has finite dimension over $k$, all the algebras $A/I_A^{\{p^n\}}$ are finite-dimensional $k$-vector spaces. One therefore sees that: \emph{every infinitesimal formal group $G$ over a field of characteristic $p>0$, such that $\omega_{G/k}$ has finite dimension over $k$, is an inductive limit of finite formal groups} (\ie of finite length, cf. 1.2.6).

\sgasection{5}{Homogeneous spaces of infinitesimal formal groups over a field}
\numpara{5.0}\nde{We have added subsection 5.0, to express the following Proposition 5.1 in the language of cocommutative Hopf algebras, and to cite the results obtained since in this direction.}Suppose, for simplicity, that $k$ is a field. Let $G$ be a formal $k$-group with affine algebra $A=\mathscr A(G)$. Let $A^+$ be the augmentation ideal of $A$; for every $a\in A$, denote by $\overline a=a-\varepsilon_A(a)1_A$ its projection onto $A^+$. Let $F$ be a formal subgroup of $G$, defined by the closed ideal $J$ of $A$, and let $\pi$ (resp. $\overline\pi$) be the projection $A\to A/J$ (resp. the composite of the projections $A\to A^+\to A^+/J$). Note that, for every $a\in A$, the projection of $\Delta_A(a)$ onto $A\widehat\otimes A^+$ is $\Delta_A(a)-a\widehat\otimes1$.

By Theorem 2.4, one can form the quotient formal $k$-variety $X=G/F$; its affine algebra $B$ is
\[
\begin{aligned}
B&=\Ker\left(\xymatrix@C=52pt{A\ar@<.6ex>[r]^{\tau_1}\ar@<-.6ex>[r]_{(\id_A\widehat\otimes\pi)\Delta_A}&A\widehat\otimes_k(A/J)}\right)\\
&=\{a\in A\mid(\id_A\widehat\otimes\pi)\Delta_A(a)=a\widehat\otimes\pi(1)\}\\
&=\Ker((\id_A\widehat\otimes\overline\pi)\Delta_A).
\end{aligned}
\]
It is also the subalgebra of $A$ consisting of the elements $\phi$ such that, for every $C\in\Ob\mathbf{Alf}/k$ and $g\in G(C)$, $h\in F(C)$, one has $\phi(gh)=\phi(g)$. For every $g,g'\in G(C)$ and $h\in F(C)$, one has $\phi(gg'h)=\phi(gg')$, so $\Delta(\phi)$ belongs to the kernel of $\id_A\widehat\otimes_k(\id_A\widehat\otimes\overline\pi)\Delta_A$, which equals $A\widehat\otimes_k B$ since $A$ is topologically flat over $k$. One therefore has $\Delta_A(B)\subset A\widehat\otimes_k B$, \ie the closed subalgebra $B$ is also a \emph{left coideal}.

On the other hand, $B$ determines $F$ since, by Corollary 2.4.1, one has $J=\overline{AB^+}$, \ie $J$ is the closed ideal generated by $B^+=B\cap A^+$. One thus obtains an \emph{injective} map $\mathcal Q$ from the set $\mathscr F$ of formal subgroups of $G$ into the set $\mathscr B$ of closed subalgebras $B$ of $A$ which are left coideals. There then arises the question of determining the image of this map, and Proposition 5.1 below shows that $\mathcal Q$ is bijective when $G$ is \emph{infinitesimal}. In fact, this is true for \emph{every} formal $k$-group $G$.

Indeed, recall (cf. 2.2.1) that the functor $G\to H(G)$ is an \emph{equivalence} between the category of formal $k$-groups and that of cocommutative $k$-Hopf algebras; if $F$ is a formal subgroup of $G$, defined by the closed ideal $J$ of $A$, then the Hopf subalgebra $H(F)$ of $H=H(G)$ is the orthogonal of $J$ for the duality between $A=H^*$ and $H=A^\dagger$ (cf. 0.2.2). On the other hand, if $B$ is a closed subalgebra of $A$ which is also a left coideal, then its orthogonal $I=B^\perp$ is a coideal of $H$ (\ie $\Delta_H(I)\subset I\otimes H+H\otimes I$ and $\varepsilon_H(I)=0$) and a left ideal. Denote by $\mathscr H$ (resp. $\mathscr I$) the set of Hopf subalgebras (resp. left ideals which are coideals) of $H$. For every $I\in\mathscr I$, denote by $\pi_I$ (resp. $\overline\pi_I$) the projection $H\to H/I$ (resp. the composite of the projections $H\to H^+\to H^+/I$), where $H^+$ is the augmentation ideal of $H$.

Let $K$ be a Hopf subalgebra of $H$ and $K^+=K\cap H^+$. If $F$ is the formal subgroup corresponding to $K$, then $J=K^\perp$ and $A^+/J$ is identified with the dual of $K^+$, and since $B=\mathcal Q(F)$ is the kernel of the map
\[
\xymatrix@C=30pt{A\ar[r]^{\Delta_A}&A\widehat\otimes_k A\ar[r]^{\id\widehat\otimes\overline\pi}&A\widehat\otimes_k(A^+/J),}
\]
one obtains that $\mathcal Q$ corresponds by duality to the map $\Phi$ which associates to $K$ the image of
\[
\xymatrix@C=30pt{H&H\otimes_k H\ar[l]_{m_H}&H\otimes_k K^+\ar[l]_{\id\widehat\otimes\mathrm{can.}},}
\]
\ie the left ideal $HK^+$, which is also a coideal. One sees similarly that the map which associates to $B\in\mathscr B$ the ideal $\overline{AB^+}$ corresponds by duality to the map $\Psi$ which associates to every $I\in\mathscr I$ the kernel of $(\id\otimes_k\overline\pi_I)\circ\Delta_H$, \ie one has
\[
\Psi(I)=\{x\in H\mid(\id\otimes_k\pi_I)\Delta_H(x)=x\otimes\pi_I(1)\}.
\]
One then has the following theorem:
\begin{theorem}{5.0.1}\label{VIIB.5.0.1}
Let $k$ be a field and $H$ a cocommutative Hopf algebra. Then the maps $\Phi$ and $\Psi$ above are inverse bijections between the set of Hopf subalgebras of $H$ and that of left ideals which are coideals.
\end{theorem}
This theorem was first proved by K. Newman, cf. \cite{Ne75}, Th. 4.1 (where the word ``cocommutative'' was omitted). His proof uses ``the Cartier--Gabriel--Kostant theorem'' (cf. 2.9) to reduce to the case where $H$ is connected, then the existence in this case of a ``Sweedler basis'' (cf. \cite{Sw67}, Th. 3), a result dual to the Dieudonné--Cartier Theorem 5.2.2 below. Another, shorter, proof was given by H. J. Schneider \cite{Sch90}, Th. 4.15. A generalization was subsequently obtained by A. Masuoka when one assumes only that the coradical $H_0$ of $H$ (\ie the sum of the simple subcoalgebras) is commutative, \cite{Ma91}, Th. 1.3 (3).

Let us mention finally that for a commutative $k$-Hopf algebra, thus corresponding to an affine group $k$-scheme $G$, one cannot expect an analogue of 5.0.1 without additional hypotheses, since for a $k$-subgroup $F$ of $G$ the quotient $G/F$ is not necessarily affine. But M. Takeuchi established in \cite{Tak72}, Th. 4.3 (resp. \cite{Tak79}, Th. 3), an analogous bijection between the set of $k$-subgroups $F$ of $G$ which are invariant (resp. such that $G/F$ is affine), and that of subalgebras $B$ of $\mathcal O(G)$ such that $\Delta(B)\subset A\otimes B$ and which are stable under the antipode (resp. and such that $B\to A$ is faithfully flat).

\numpara{5.1}Let $k$ be a pseudocompact ring.\nde{In the original, $k$ is assumed in 5.1 to be a field. In fact, this hypothesis can be replaced by flatness hypotheses; we have modified nos. 5.1 to 5.1.5 accordingly.} Let $G$ be a topologically flat infinitesimal formal $k$-group, $A$ its affine algebra, $B$ a closed subalgebra of $A$, $X=\Spf(B)$ and $r:G\to X$ the epimorphism induced by the inclusion of $B$ into $A$. One proposes to see under what condition $r$ makes $X$ the right quotient of $G$ by a subgroup $H$ (cf. 2.4).\nde{We have replaced ``left'' by ``right'' and modified the statement of Proposition 5.1, to bring out more clearly, on the one hand, the equivalent conditions (i), (ii), and, on the other hand, the conclusion $\Spf(B)\simeq G/H$.}
\begin{proposition}{}\label{VIIB.5.1}
Let $G$ be a topologically flat infinitesimal formal $k$-group, $A$ its affine algebra, $I_A$ the augmentation ideal of $A$, $B$ a closed subalgebra of $A$, and $J_B=\overline{AI_B}$, where $I_B=B\cap I_A$. Suppose that $A$ is topologically flat over $k$, as is $\overline{J_B^n}/\overline{J_B^{n+1}}$ for every $n\ge0$. Then the following two assertions are equivalent:
\begin{enumeratei}
\item For every $x\in B$, $\Delta_A(x)-x\widehat\otimes1$ belongs to $A\widehat\otimes_k J_B$.
